% ===== Section 4: Why your ears are burning (attention, ~18 min) =====
\topicopener{headphones}{Ուշադրություն}{Ինչու են ականջներդ «այրվում»}{Ուշադրության ուժն ու թուլությունները, ու ինչու չես կարող չլսել ուրիշի խոսակցությունը}{thinking}

\begin{frame}{Ուղեղը չի կարող ամեն ինչ մշակել}
  \begin{columns}[T]
    \begin{column}{0.52\textwidth}
      \begin{itemize}
        \item Զգայարանները անընդհատ \textbf{հսկայական} տեղեկություն են ուղարկում
        \item<2-> Ուղեղը շատ «թանկ» օրգան է. չորացող ներկին նայելու վրա ռեսուրս վատնելը անիմաստ է
        \item<3-> \textbf{Ուշադրությունը} որոշում է, թե ինչը կնկատենք։ Ու, որ հաճախ ավելի կարևոր է, ինչը \textbf{չենք նկատի}
      \end{itemize}
    \end{column}
    \begin{column}{0.44\textwidth}
      \onslide<4->{%
      \begin{card}[colback=butter!45!cream]
        {\blackfont Երկու հարց}\par\smallskip
        \emoji{battery}\ \textbf{Որքա՞ն} ենք կարողանում մշակել\par\smallskip
        \emoji{target}\ \textbf{Ի՞նչն է} որոշում, թե ուր ուղղվի ուշադրությունը
      \end{card}}
    \end{column}
  \end{columns}
\end{frame}

\begin{frame}{Բոլորը միանգամից խոսում են}
  \centering
  % Cherry's dichotic listening: a different stream of words in each ear; repeat the left one.
% What people remember from the right ear appears on slide 2 (use inside a frame).
\begin{tikzpicture}
  \path[use as bounding box] (-6.4,-1.65) rectangle (6.4,1.8);
  \node at (0,0.05) {\mascot[2.3cm]{thinking}};
  % headphones
  \draw[plum, line width=2.2pt] (-1.15,0.15) .. controls (-1.1,1.75) and (1.1,1.75) .. (1.15,0.15);
  \filldraw[fill=amBlue, draw=plum, line width=1pt, rounded corners=1.5mm] (-1.42,-0.45) rectangle (-0.98,0.4);
  \filldraw[fill=plum!45!cream, draw=plum, line width=1pt, rounded corners=1.5mm] (0.98,-0.45) rectangle (1.42,0.4);
  % left ear: the stream to repeat
  \node[chatmsg=3.7cm, fill=amBlue!12!cardwhite, align=center] (L) at (-4.2,0.55) {«խնձոր... սեղան... ամպ... գետ...»};
  \node[wpill=amBlue] at (-4.2,1.5) {ձախ ականջ. կրկնիր բարձրաձայն};
  \draw[warr=amBlue] (L.east) -- (-1.5,0.1);
  % right ear: the ignored stream
  \node[chatmsg=3.7cm, fill=cardwhite, align=center, text=plum!60!cream] (R) at (4.2,0.55) {«շուն... լամպ... կոշիկ... ձյուն...»};
  \node[wpill=plum!60] at (4.2,1.5) {աջ ականջ. անտեսիր};
  \draw[warr=plum!50] (R.west) -- (1.5,0.1);
  % what is left of the ignored stream
  \node[visible on=<2->, draw=pinkdark, line width=1.2pt, fill=blush!15!cardwhite, rounded corners=2.5mm,
        text width=4.0cm, align=center, font=\small, text=plum, inner sep=4pt] at (4.2,-0.95)
    {Աջից հիշում են միայն՝ \textbf{տղամարդ էր, թե կին}։ Նույնիսկ \textbf{լեզուն} չէին նկատում};
\end{tikzpicture}
\par\smallskip
  \begin{columns}[T]
    \begin{column}{0.47\textwidth}
      \small 1953. Քոլին Չերին (Colin Cherry) մարդկանց երկու ականջին տարբեր բառեր էր տալիս (dichotic listening)։
    \end{column}
    \begin{column}{0.47\textwidth}
      \onslide<3->{\small Եզրակացություն՝ ուշադրությունը \textbf{նեղ բկով շիշ} է (bottleneck)։ Հեռադիտակի պես. մի փոքր կտորը մանրամասն, մնացածը՝ ոչինչ։}
    \end{column}
  \end{columns}
\end{frame}

\begin{frame}{Բայց ուղեղը լսում էր}
  \begin{columns}[T]
    \begin{column}{0.52\textwidth}
      \begin{itemize}
        \item 1975. Ֆոն Ռայթն (von Wright) ու գործընկերները մարդկանց սովորեցրին որոշ բառերից հետո \textbf{էլեկտրական հարված} սպասել \emoji{lightning}
        \item<2-> Հետո նույն բառերը հնչեցին \textbf{անտեսվող} ականջում
        \item<3-> Մարդիկ բառերը չէին նկատում, բայց մարմինը \textbf{վախով} էր արձագանքում
      \end{itemize}
    \end{column}
    \begin{column}{0.44\textwidth}
      \onslide<4->{%
      \begin{card}[colback=brainpink!25!cardwhite]
        Ուրեմն ուղեղը մշակում էր նաև «մյուս» ականջը, պարզապես տեղեկությունը \textbf{գիտակցությանը չէր հասնում}։ Շշի բկի մոդելը չի աշխատում։
      \end{card}}
    \end{column}
  \end{columns}
\end{frame}

\begin{frame}{Կոկտեյլի երեկույթի էֆեկտ}
  \begin{columns}[T]
    \begin{column}{0.5\textwidth}
      \begin{poll}
        Բարձրացրու ձեռքդ, եթե աղմկոտ սենյակում \textbf{լսել ես անունդ} ուրիշների խոսակցությունից ու միանգամից «միացել» նրանց զրույցին։
      \end{poll}
      \vspace{1mm}
      \onslide<3->{\small Սա կոչվում է \textbf{կոկտեյլի երեկույթի էֆեկտ} (cocktail party effect). «որովհետև հոգեբանները նրբաճաշակ մարդիկ են»։}
    \end{column}
    \begin{column}{0.46\textwidth}
      \onslide<2->{%
      \begin{card}
        \emoji{party}\ Հարսանիք։ Սեղանի մոտ խոսում ես ֆուտբոլից կամ նեխուրից։ Հարևան խմբում ինչ-որ մեկն ասում է \textbf{քո անունը}, ու հետո՝ «...իսկական մաշկի վատնում է»։\par\smallskip
        Ու հիմա դու լսում ես միայն \textbf{նրանց} ու մտածում, թե ինչու ես այդ մարդուն քավոր ընտրել։
      \end{card}}
    \end{column}
  \end{columns}
\end{frame}

\begin{frame}{Հարսանիքում}
  \centering
  \begin{meme}[6.9cm]{monkey_puppet.jpg}
    \memecap[text width=5.8cm, font=\blackfont\normalsize]{0.5}{0.84}{Հարևան սեղանից ինչ-որ մեկը ասում է իմ անունը}
    \memecap{0.27}{0.08}{ես, իբր չեմ լսում}
  \end{meme}
\end{frame}

\begin{frame}{Ուշադրությունը՝ սահմանափակ բյուջե}
  \centering
  % Kahneman's capacity model: attention as a limited budget shared between tasks. Each task takes a
% share; the fourth one does not fit. Segments appear on slides 1-4 (use inside a frame).
\begin{tikzpicture}[seg/.style={draw=plum, line width=0.8pt, font=\small, text=plum, align=center, minimum height=0.9cm}]
  \path[use as bounding box] (-0.2,-1.0) rectangle (12.4,1.2);
  \draw[plum, line width=1.4pt, rounded corners=2mm] (0,-0.5) rectangle (9.6,0.5);
  \node[font=\blackfont\scriptsize, text=plum, anchor=south west] at (0,0.55) {ուշադրության «բյուջեն»};
  \node[seg, fill=mint!60!cardwhite, minimum width=4.2cm, anchor=west] at (0.05,0) {\emoji{car}\ ճանապարհ};
  \node[seg, fill=butter!70!cardwhite, minimum width=3.0cm, anchor=west, visible on=<2->] at (4.25,0) {\emoji{speech}\ զրույց};
  \node[seg, fill=peach!60!cardwhite, minimum width=2.25cm, anchor=west, visible on=<3->] at (7.25,0) {\emoji{musical}\ ռադիո};
  \node[seg, fill=blush!50!cardwhite, draw=pinkdark, minimum width=2.6cm, anchor=west, visible on=<4->] at (9.7,0) {\emoji{phone}\ հեռախոս};
  \node[visible on=<4->, font=\blackfont\scriptsize, text=pinkdark] at (11.0,-0.8) {չի տեղավորվում};
\end{tikzpicture}
\par\smallskip
  \begin{columns}[T]
    \begin{column}{0.47\textwidth}
      \begin{card}[colback=lavender!35!cardwhite]
        \small\textbf{Ռեսուրսի մոդել} (Kahneman, 1973). ուշադրությունը կարելի է \textbf{բաժանել} մի քանի բանի միջև, քանի դեռ ռեսուրսը չի սպառվել։
      \end{card}
    \end{column}
    \begin{column}{0.47\textwidth}
      \onslide<5->{%
      \begin{card}[colback=mint!35!cardwhite]
        \small Ինչու՞ է այդքան քիչ։ Ուշադրությունը տեղեկություն է տալիս \textbf{կարճաժամկետ հիշողությանը}, իսկ այն տեղավորում է մոտ 4 բան (վեբինար 2)։
      \end{card}}
    \end{column}
  \end{columns}
\end{frame}

\begin{frame}{Հեռախոսով, բայց առանց ձեռքերի}
  \begin{columns}[c]
    \begin{column}{0.5\textwidth}
      \begin{itemize}
        \item Շատ երկրներում վարելիս հեռախոսը ձեռքում պահել չի կարելի, իսկ առանց ձեռքերի (hands-free)՝ կարելի է
        \item<2-> Յուտայի համալսարան, 2013. hands-free-ն \textbf{նույնքան վատ է}. արձագանքը, շուրջը նայելը, նշանները նկատելը նույնքան են տուժում
        \item<3-> Ըստ Բըրնեթի՝ ռադիոն փոխելն ու ուղևորի հետ զրույցն էլ են նույնքան շեղում
      \end{itemize}
    \end{column}
    \begin{column}{0.46\textwidth}
      \centering
      \onslide<4->{%
      \begin{meme}[5.6cm]{left_exit.jpg}
        \memecap{0.45}{0.83}{ճանապարհ}
        \memecap{0.62}{0.69}{նոր հաղորդագրություն}
        \memecap{0.48}{0.12}{ուշադրությունս}
      \end{meme}}
    \end{column}
  \end{columns}
\end{frame}

\begin{frame}{Ինչպե՞ս ենք ընդհանրապես վարում}
  \begin{columns}[c]
    \begin{column}{0.58\textwidth}
      \begin{itemize}
        \item Սկսնակի համար վարելը սարսափելի ծանր է
        \item<2-> Բայց կրկնությունից հետո այն դառնում է \textbf{պրոցեդուրային հիշողություն} (վեբինար 2), ու ենթագիտակցությունն է վարում «առանց մտածելու»
        \item<3-> Բայց ոչ ամբողջովին. մյուս վարորդներն ու վտանգները ամեն անգամ նոր են, ու դրանց համար \textbf{գիտակցված ուշադրություն} է պետք
      \end{itemize}
    \end{column}
    \begin{column}{0.38\textwidth}
      \centering
      \bigemoji[1.6cm]{car}\par\medskip
      \onslide<3->{%
      \begin{card}[colback=peach!40!cardwhite]
        \small Աչքերդ կարող են ճանապարհին լինել, բայց եթե չես «նայում» նրանց տվածին, \textbf{անիմաստ} է։
      \end{card}}
    \end{column}
  \end{columns}
\end{frame}

\begin{frame}{Որտե՞ղ է ուշադրությունը}
  \begin{columns}[c]
    \begin{column}{0.48\textwidth}
      \centering
      \begin{tikzpicture}
        \path[use as bounding box] (-3.1,-1.8) rectangle (3.1,1.9);
        \begin{scope}[scale=0.56]
          \brainlobes{lavender!40!cream}{mint!40!cream}{peach!45!cream}{butter!80!cream}
          \fill[amBlue, opacity=0.55] (-2.3,1.0) ellipse (0.8 and 0.65);
          \draw[amOrange!85!black, line width=1.4pt, dashed, visible on=<2->] (-0.9,0.9) .. controls (-0.2,1.45) and (0.8,1.35) .. (1.3,0.95);
          \node[wpill=amBlue] at (-2.3,2.75) {նախաճակատային կեղև};
          \node[wpill=amOrange!85!black, visible on=<2->] at (1.7,-2.95) {առաջային գոտկային կեղև (ներսում)};
        \end{scope}
      \end{tikzpicture}
    \end{column}
    \begin{column}{0.48\textwidth}
      \begin{itemize}
        \item \textbf{Նախաճակատային կեղև}. այնտեղ է կարճաժամկետ հիշողությունը
        \item<2-> \textbf{Առաջային գոտկային կեղև} (anterior cingulate)
      \end{itemize}
      \onslide<3->{\small Համակարգերը \textbf{ցրված} են ուղեղով մեկ։\par\smallskip}
      \onslide<4->{%
      \begin{scienceupdate}
        Գրքում՝ «քունքային բլթի խորքում»։ Իրականում այն կիսագնդերի \textbf{ներքին} մակերեսին է։
      \end{scienceupdate}}
    \end{column}
  \end{columns}
  \onslide<4->{\source{Bush, Luu \& Posner, Trends Cogn. Sci., 2000}}
\end{frame}

\begin{frame}{Ներքևից վեր և վերևից ներքև}
  \begin{columns}[T]
    \begin{column}{0.47\textwidth}
      \begin{card}[colback=peach!40!cardwhite]
        \emoji[1.3em]{bell}\ {\blackfont Ներքևից վեր} (bottom-up)\par\smallskip
        Դրսի ինչ-որ բան \textbf{ինքն է} գրավում ուշադրությունը։ Անունդ հնչեց, ու դու արդեն լսում ես։
      \end{card}
    \end{column}
    \begin{column}{0.47\textwidth}
      \onslide<2->{%
      \begin{card}[colback=lavender!40!cardwhite]
        \emoji[1.3em]{brain}\ {\blackfont Վերևից ներքև} (top-down)\par\smallskip
        \textbf{Դու ես} որոշում, թե ուր նայել ու ինչ լսել։ Ուզում ես իմանալ, թե ինչ են ասում քո մասին։
      \end{card}}
    \end{column}
  \end{columns}
  \vspace{2mm}
  \onslide<3->{%
  \begin{card}[colback=mint!35!cardwhite]
    \small\emoji{headphones}\ Ինչպե՞ս ենք «ուղղում» ականջները։ 2012-ին Չանգն ու Մեսգարանին երեք պացիենտի լսողական կեղևում տեսան, որ ակտիվություն է տալիս միայն այն ձայնը, որին մարդը \textbf{ուշադիր է}։ Ուղեղն իսկապես կարող է «անջատել» ձանձրալի զրուցակցին՝ ոզնիներ դիտելու իր հոբբիով։
  \end{card}}
  \onslide<3->{\source{Mesgarani \& Chang, Nature, 2012}}
\end{frame}

\begin{frame}{Աչքերը՝ ուշադրության գլխավոր գործիքը}
  \begin{columns}[T]
    \begin{column}{0.47\textwidth}
      \begin{card}[colback=lavender!40!cardwhite]
        \emoji[1.3em]{target}\ {\blackfont Նպատակով}\par\smallskip
        Տեսար «ԱՆՎՃԱՐ ԲԵԿՈՆ» ցուցանակը, ու աչքերդ \textbf{դու ես} ուղղում այնտեղ։
      \end{card}
    \end{column}
    \begin{column}{0.47\textwidth}
      \onslide<2->{%
      \begin{card}[colback=peach!40!cardwhite]
        \emoji[1.3em]{tiger}\ {\blackfont Ինքնաբերաբար}\par\smallskip
        Վագրի մռնչյուն, ճյուղի ճռռոց. ուշադրությունը շրջվում է \textbf{նախքան} գիտակցությունը կհասկանա։
      \end{card}}
    \end{column}
  \end{columns}
  \vspace{2mm}
  \onslide<3->{%
  \begin{scienceupdate}[width=0.86\textwidth]
    Գրքում սրանք կոչվում են «overt» ու «covert»։ Այսօր overt-ը նշանակում է \textbf{աչքերը շարժելով}, covert-ը՝ առանց շարժելու։
  \end{scienceupdate}}
  \onslide<3->{\source{Posner, Q. J. Exp. Psychol., 1980}}
\end{frame}

\begin{frame}{Ծնողը, հեռուստացույցը ու աղբը}
  \centering
  % Burnett's parent-and-bins version of how attention shifts to something important: the posterior
% parietal cortex disengages, the superior colliculus moves, the pulvinar re-engages. Steps on slides 1-3.
\begin{tikzpicture}[
    stage/.style={draw=plum, line width=0.9pt, rounded corners=3mm, text width=3.6cm, align=center,
      font=\small, text=plum, inner sep=3.5pt, minimum height=2.7cm},
  ]
  \node[stage, fill=lavender!40!cardwhite] (a) at (0,0)
    {\bigemoji[0.7cm]{tv}\\[1mm] \textbf{Անջատել}\\ {\scriptsize հետին գագաթային կեղև}\\[1mm] Ծնողն անջատում է հեռուստացույցը};
  \node[stage, fill=mint!40!cardwhite, visible on=<2->] (b) at (4.6,0)
    {\bigemoji[0.7cm]{walking}\\[1mm] \textbf{Տեղափոխել}\\ {\scriptsize վերին թմբիկներ (superior colliculus)}\\[1mm] Երեխային տանում է խոհանոց};
  \node[stage, fill=peach!45!cardwhite, visible on=<3->] (c) at (9.2,0)
    {\bigemoji[0.7cm]{wastebasket}\\[1mm] \textbf{Նորից միացնել}\\ {\scriptsize թալամուսի բարձիկ (pulvinar)}\\[1mm] Աղբի տոպրակը տալիս է ձեռքը ու հրում դեպի դուռը};
  \draw[warr, visible on=<2->] (a) -- (b);
  \draw[warr, visible on=<3->] (b) -- (c);
\end{tikzpicture}
\par\smallskip
  \onslide<4->{%
  \begin{card}[width=0.9\textwidth, colback=butter!40!cream]
    \centering Այս համակարգը կարող է \textbf{վերագրել} գիտակցականին։ Անծանոթ ուրվագիծը կարող է հարձակվող լինել։ Կամ ձանձրալի գործընկերը, որ պատմում է իր ոտքի սնկի մասին։
  \end{card}}
\end{frame}

\begin{frame}{Աչքի անկյունում՝ սարդ}
  \begin{columns}[c]
    \begin{column}{0.56\textwidth}
      \begin{itemize}
        \item Աշխատում ես, ու աչքիդ անկյունում ինչ-որ բան շարժվեց՝ մեծ \textbf{սարդի} չափ
        \item<2-> Չես նայում (հանկարծ իսկապես սարդ է), բայց շարունակում ես մտքով «հետևել» այդ անկյունին
        \item<3-> Ուշադրությունը \textbf{կապված չէ} աչքերի ուղղությանը
      \end{itemize}
    \end{column}
    \begin{column}{0.4\textwidth}
      \centering
      \bigemoji[1.2cm]{computer}\ \onslide<1->{\raisebox{-2mm}{\bigemoji[0.7cm]{spider}}}\par\medskip
      \onslide<4->{%
      \begin{card}[colback=lavender!35!cardwhite]
        \emoji{firework}\ Թե ինչն է «կարևոր», որոշում է համատեքստը։ Ամանորի գիշերը հրավառության ձայնը սովորական է. ավելի տարօրինակ կլիներ \textbf{լռությունը}։
      \end{card}}
    \end{column}
  \end{columns}
\end{frame}

\begin{frame}{Երկու բան միաժամանակ}
  \begin{columns}[c]
    \begin{column}{0.5\textwidth}
      \centering
      % A Posner cueing trial: look at the cross; an arrow cue points left; the target appears on the right
% (a wrong cue), so attention has to move. Slides: 1 = boxes, 2 = cue, 3 = target (use inside a frame).
\begin{tikzpicture}[scale=0.8]
  \path[use as bounding box] (-4.1,-1.3) rectangle (4.1,1.3);
  \filldraw[fill=cardwhite, draw=plum, line width=1pt, rounded corners=1mm] (-3.9,-0.9) rectangle (-2.1,0.9);
  \filldraw[fill=cardwhite, draw=plum, line width=1pt, rounded corners=1mm] (2.1,-0.9) rectangle (3.9,0.9);
  \draw[plum, line width=1.5pt] (-0.25,0) -- (0.25,0) (0,-0.25) -- (0,0.25);
  \draw[amBlue, line width=3pt, -{Stealth[length=4mm]}, visible on=<2->] (0.9,0.6) -- (-0.9,0.6);
  \node[visible on=<3->] at (3.0,0) {\bigemoji[1.1cm]{star}};
\end{tikzpicture}
\par\smallskip
      {\small Նայիր \textbf{+}-ին։ Սլաքը ցույց է տալիս, թե որտեղ կհայտնվի աստղը... երբեմն սխալ}
    \end{column}
    \begin{column}{0.46\textwidth}
      \begin{itemize}
        \item<4-> Մայքլ Պոզների (Michael Posner) թեստերում արդեն \textbf{երկու} հուշումը շատ է
        \item<5-> Տեսողական ու լսողական առաջադրանքները կարելի է բաժանել, բայց միայն պարզ «այո/ոչ»-ի դեպքում
        \item<6-> Բացառություն՝ երբ մեկը շատ վարժ է. փորձառու վարորդը զրուցում է, փորձառու մեքենագրողը թվաբանություն է անում
      \end{itemize}
    \end{column}
  \end{columns}
\end{frame}

\begin{frame}{Օձեր ու սարդեր. վայրկյանի մի մասում}
  \begin{columns}[T]
    \begin{column}{0.52\textwidth}
      \begin{itemize}
        \item Ուփսալայի համալսարան. մարդկանց ցույց էին տալիս օձերի ու սարդերի նկարներ \textbf{1/300 վայրկյանից} էլ կարճ
        \item<2-> Գիտակցորեն ճանաչելու համար պետք է մոտ \textbf{կես վայրկյան}
        \item<3-> Բայց ափերը միևնույն է \textbf{քրտնում} էին. անգիտակից ուշադրությունը վտանգը նկատում է նախքան գիտակցությունը կասի՝ «Հը՞, ի՞նչ»
      \end{itemize}
    \end{column}
    \begin{column}{0.44\textwidth}
      \centering
      \bigemoji[1.1cm]{snake}\ \bigemoji[1.1cm]{spider}\par\medskip
      \onslide<4->{%
      \begin{scienceupdate}
        Այս տիպի դասական փորձերում նկարը ցույց էին տալիս \textbf{30 միլիվայրկյան} (մոտ 1/33 վ), հետո՝ ծածկող պատկեր։ «1/300»-ը չկարողացանք հաստատել։
      \end{scienceupdate}}
    \end{column}
  \end{columns}
  \onslide<4->{\source{Öhman \& Soares, J. Abnorm. Psychol., 1994}}
\end{frame}

\begin{frame}{Դռան փորձը}
  \centering
  % Simons & Levin's door study in three panels: an experimenter asks a pedestrian for directions,
% two people carry a door between them, and a different experimenter is standing there afterwards.
% Panels on slides 1-3 (use inside a frame).
\begin{tikzpicture}[scale=0.84,
    panel/.style={draw=plum, line width=0.8pt, fill=cardwhite, rounded corners=2.5mm},
    plbl/.style={font=\scriptsize, text=plum, align=center, text width=3.1cm},
  ]
  \path[use as bounding box] (-0.1,-0.65) rectangle (12.6,2.85);
  \foreach \i/\s in {0/1, 1/2, 2/3}{
    \begin{scope}[xshift=\i*4.25cm, visible on=<\s->]
      \draw[panel] (0,0) rectangle (4.0,2.9);
      % pedestrian (right) with the map
      \fill[plum!70] (3.0,1.95) circle (0.24);
      \fill[plum!70, rounded corners=1.5mm] (2.72,0.35) rectangle (3.28,1.65);
      \filldraw[fill=butter, draw=plum, line width=0.6pt, rotate around={-12:(2.45,1.15)}] (2.15,0.9) rectangle (2.75,1.4);
    \end{scope}}
  % panel 1: experimenter A (blue)
  \fill[amBlue] (1.0,1.95) circle (0.24);
  \fill[plum] (0.78,2.05) rectangle (1.22,2.2);
  \fill[amBlue, rounded corners=1.5mm] (0.72,0.35) rectangle (1.28,1.65);
  \node[say, font=\scriptsize, callout absolute pointer={(1.2,2.0)}] at (1.9,2.55) {Ինչպե՞ս գնամ...};
  \node[plbl] at (2.0,-0.35) {հարցնում է ճանապարհը};
  % panel 2: the door
  \begin{scope}[xshift=4.25cm, visible on=<2->]
    \fill[amBlue] (1.0,1.95) circle (0.24);
    \fill[amBlue, rounded corners=1.5mm] (0.72,0.35) rectangle (1.28,1.65);
    \filldraw[fill=peachdark!70, draw=plum, line width=0.9pt] (1.55,0.2) rectangle (2.35,2.55);
    \fill[butter] (2.2,1.35) circle (0.06);
    \node[plbl] at (2.0,-0.35) {մեջտեղով դուռ են անցկացնում};
  \end{scope}
  % panel 3: experimenter B (orange, different hair)
  \begin{scope}[xshift=8.5cm, visible on=<3->]
    \fill[amOrange] (1.0,1.95) circle (0.24);
    \fill[peachdark] (0.74,2.1) .. controls (1.0,2.4) .. (1.26,2.1) -- cycle;
    \fill[amOrange, rounded corners=1.5mm] (0.68,0.35) rectangle (1.32,1.75);
    \node[say, font=\scriptsize, callout absolute pointer={(1.2,2.0)}] at (1.9,2.55) {...ու հետո՞};
    \node[plbl] at (2.0,-0.35) {ուրիշ մարդ է կանգնած};
  \end{scope}
\end{tikzpicture}
\par
  \begin{columns}[T]
    \begin{column}{0.47\textwidth}
      \begin{poll}
        Կնկատեի՞ր, որ դիմացդ \textbf{ուրիշ մարդ} է։ Բարձրացրու ձեռքդ, եթե այո։
      \end{poll}
    \end{column}
    \begin{column}{0.47\textwidth}
      \onslide<4->{%
      \begin{card}[colback=brainpink!25!cardwhite]
        \small 1998. Նկատել է \textbf{15-ից միայն 7-ը}։ Իսկ 50 ուսանողից բոլորն էին վստահ, որ կնկատեն։
      \end{card}}
    \end{column}
  \end{columns}
  \onslide<4->{\source{Simons \& Levin, Psychon. Bull. Rev., 1998}}
\end{frame}

\begin{frame}{«Դռան փորձ»}
  \begin{columns}[c]
    \begin{column}{0.56\textwidth}
      \begin{bookquote}
        Այն կոչվում է «դռան փորձ», որովհետև, ըստ երևույթին, ամենահետաքրքիրը դուռն է։ Գիտնականները տարօրինակ մարդիկ են։
      \end{bookquote}
    \end{column}
    \begin{column}{0.4\textwidth}
      \centering
      \bigemoji[1.8cm]{door}\par\medskip
      {\small Ու մի մանրուք հոդվածից. մարդկանց մեծ մասն այդ պահից նկատել էր միայն, որ դուռ տանողները \textbf{անքաղաքավարի} էին։}
    \end{column}
  \end{columns}
  \source{Simons \& Levin, 1998}
\end{frame}

\begin{frame}{Փոփոխության կուրություն}
  \begin{columns}[c]
    \begin{column}{0.66\textwidth}
      \centering
      % Change-blindness demo (flicker): scene with a chimney on slides 1, 5 and 8, without it on 3 and 7,
% a blank screen on 2, 4 and 6; slide 8 circles the change. Use inside a frame with 8 slides.
\begin{tikzpicture}[scale=0.9]
  \path[use as bounding box] (-5.2,-2.15) rectangle (5.2,2.15);
  \begin{scope}
    \clip[rounded corners=3mm] (-5.1,-2.1) rectangle (5.1,2.1);
    \fill[lavender!40!cardwhite] (-5.1,-2.1) rectangle (5.1,2.1);
    \fill[mint!60!cardwhite] (-5.1,-2.1) rectangle (5.1,-0.9);
    % sun and cloud
    \fill[amOrange!75] (3.9,1.45) circle (0.4);
    \fill[cardwhite] (-3.2,1.45) ellipse (0.75 and 0.3) (-2.7,1.6) ellipse (0.5 and 0.3);
    % house
    \only<1,2,4,5,6,8>{\filldraw[fill=peachdark!75, draw=plum, line width=0.8pt] (0.95,0.95) rectangle (1.3,1.65);}
    \filldraw[fill=peach, draw=plum, line width=0.8pt] (-0.4,-1.0) rectangle (1.9,0.55);
    \filldraw[fill=plum!70, draw=plum, line width=0.8pt] (-0.6,0.55) -- (0.75,1.55) -- (2.1,0.55) -- cycle;
    \filldraw[fill=butter, draw=plum, line width=0.7pt] (-0.1,-0.35) rectangle (0.45,0.2);
    \filldraw[fill=peachdark, draw=plum, line width=0.7pt] (0.95,-1.0) rectangle (1.45,-0.05);
    % tree, fence, car
    \fill[peachdark] (-2.75,-1.0) rectangle (-2.5,-0.2);
    \fill[mintdark] (-2.62,0.15) circle (0.62);
    \foreach \x in {-4.8,-4.4,...,-3.2}{\filldraw[fill=cardwhite, draw=plum, line width=0.5pt] (\x,-1.0) rectangle ++(0.18,0.6);}
    \draw[plum, line width=0.6pt] (-4.85,-0.65) -- (-3.0,-0.65);
    \node at (3.5,-1.15) {\bigemoji[0.9cm]{car}};
  \end{scope}
  \draw[plum, line width=0.9pt, rounded corners=3mm] (-5.1,-2.1) rectangle (5.1,2.1);
  % blank flashes
  \only<2,4,6>{\fill[cream] (-5.2,-2.15) rectangle (5.2,2.15);}
  % the answer
  \only<8>{\draw[amRed, line width=2pt] (1.12,1.3) circle (0.55);
    \node[wpill=amRed] at (2.6,1.75) {ծխնելույզը};}
\end{tikzpicture}

    \end{column}
    \begin{column}{0.3\textwidth}
      \begin{chatbox}
        Նկարը մի պահ անհետանում է։ Ի՞նչ է \textbf{փոխվել}։ Գրիր չաթում։
      \end{chatbox}
      \vspace{1mm}
      \onslide<8>{\small Ընդհատումից հետո ուղեղը տեսարանը նախորդի հետ \textbf{չի համեմատում} (change blindness)։}
    \end{column}
  \end{columns}
\end{frame}

\begin{frame}{Գլխի վերևում ցուցադրիչ}
  \begin{columns}[c]
    \begin{column}{0.56\textwidth}
      \centering
      % A cluttered heads-up display over a runway: the plane standing on the runway is easy to miss.
% Slide 2 circles it (use inside a frame).
\begin{tikzpicture}[hud/.style={text=mintdark, font=\blackfont\scriptsize, inner sep=1pt}]
  \path[use as bounding box] (-4.0,-2.1) rectangle (4.0,2.1);
  \begin{scope}
    \clip[rounded corners=4mm] (-3.9,-2.0) rectangle (3.9,2.0);
    \shade[top color=lavender!70!cardwhite, bottom color=cardwhite] (-3.9,0.1) rectangle (3.9,2.0);
    \fill[mint!45!cardwhite] (-3.9,-2.0) rectangle (3.9,0.1);
    % runway in perspective
    \fill[plum!35!cardwhite] (-1.7,-2.0) -- (1.7,-2.0) -- (0.22,0.1) -- (-0.22,0.1) -- cycle;
    \foreach \a/\b in {-1.9/-1.5, -1.1/-0.8, -0.5/-0.3, -0.12/0.0}{\fill[cardwhite] (-0.03,\a) rectangle (0.03,\b);}
    % a small plane standing on the runway
    \begin{scope}[shift={(0.18,-0.55)}, scale=0.32]
      \fill[plum] (-0.9,0) -- (0.9,0) -- (0.9,0.18) -- (-0.9,0.18) -- cycle;
      \fill[plum] (-0.15,0.05) -- (0.15,0.05) -- (0.6,-0.4) -- (0.35,-0.4) -- cycle;
      \fill[plum] (-0.15,0.05) -- (0.15,0.05) -- (-0.6,-0.4) -- (-0.35,-0.4) -- cycle;
      \fill[plum] (0.7,0.18) -- (0.9,0.18) -- (0.9,0.6) -- cycle;
    \end{scope}
    % HUD symbology
    \foreach \y/\n in {1.15/10, 0.65/5, -0.55/-5, -1.05/-10}{
      \draw[mintdark, line width=1pt] (-1.5,\y) -- (-0.55,\y) (0.55,\y) -- (1.5,\y);
      \node[hud, anchor=east] at (-1.55,\y) {\n};
      \node[hud, anchor=west] at (1.55,\y) {\n};}
    \draw[mintdark, line width=1.2pt] (-1.9,0.1) -- (-0.4,0.1) (0.4,0.1) -- (1.9,0.1);
    \draw[mintdark, line width=1.2pt] (0,-0.35) circle (0.17) (-0.45,-0.35) -- (-0.17,-0.35) (0.17,-0.35) -- (0.45,-0.35) (0,-0.18) -- (0,-0.02);
    \draw[mintdark, line width=1pt] (-3.55,-0.25) rectangle (-2.65,0.25);
    \node[hud, font=\blackfont\small] at (-3.1,0) {142};
    \draw[mintdark, line width=1pt] (2.55,-0.25) rectangle (3.6,0.25);
    \node[hud, font=\blackfont\small] at (3.07,0) {1200};
    \foreach \x/\h in {-1.6/33, -0.8/34, 0/35, 0.8/36, 1.6/01}{
      \draw[mintdark, line width=0.8pt] (\x,1.42) -- (\x,1.55);
      \node[hud, anchor=south] at (\x,1.55) {\h};}
    \node[hud, anchor=west] at (-3.75,1.6) {G 1.0};
    \node[hud, anchor=west] at (-3.75,1.3) {M 0.21};
    \node[hud, anchor=west] at (-3.75,-0.85) {ILS 110.3};
    \node[hud, anchor=west] at (-3.75,-1.15) {GS 3.0};
    \node[hud, anchor=west] at (-3.75,-1.45) {FLAPS 30};
    \node[hud, anchor=east] at (3.75,1.6) {WIND 270/12};
    \node[hud, anchor=east] at (3.75,-0.85) {BARO 1013};
    \node[hud, anchor=east] at (3.75,-1.15) {RA 210};
    \node[hud, anchor=east] at (3.75,-1.45) {V/S -700};
    \node[hud] at (0,-1.75) {GEAR DOWN};
  \end{scope}
  \draw[plum, line width=1.2pt, rounded corners=4mm] (-3.9,-2.0) rectangle (3.9,2.0);
  \begin{scope}[visible on=<2->]
    \draw[amRed, line width=2pt] (0.18,-0.6) circle (0.42);
    \node[wpill=amRed] at (1.65,-1.0) {ինքնաթիռ թռիչքուղու վրա};
  \end{scope}
\end{tikzpicture}

    \end{column}
    \begin{column}{0.4\textwidth}
      \begin{itemize}
        \item Ցուցմունքներն օդաչուի ապակու վրա են (heads-up display), ներքև նայել պետք չէ։ Ավելի ապահո՞վ է
        \item<2-> Մի քիչ ծանրաբեռնված ցուցադրիչը \textbf{սպառում է} ուշադրությունը
        \item<3-> Սիմուլյատորում վայրէջք են կատարել \textbf{ուրիշ ինքնաթիռի վրա}։ ՆԱՍԱ-ն սրա վրա հարյուրավոր միլիոններ է ծախսել
      \end{itemize}
    \end{column}
  \end{columns}
\end{frame}

\begin{frame}{Գլխի վերջը}
  \begin{columns}[c]
    \begin{column}{0.6\textwidth}
      \begin{bookquote}
        Կարող ես վիճել, որ ուշադրությունդ այդքան էլ սահմանափակ չէ։ Բայց եթե վիճում ես, ուրեմն ուշադիր չէիր։ Բարեբախտաբար, արդեն պարզեցինք, որ դրա համար քեզ չի կարելի մեղադրել։
      \end{bookquote}
    \end{column}
    \begin{column}{0.36\textwidth}
      \centering
      \mascot[3cm]{smug}
    \end{column}
  \end{columns}
\end{frame}
